\chapter{球座標系}
    \section{計量係数}
        \begin{shadebox}
            \[h_r = 1, \hth = r, \hph = r\sinth\]
        \end{shadebox}
        \begin{proof}
            \quad\par
            位置ベクトルは
            \[\vr = \dthv{r\sinth\cosph}{r\sinth\sinph}{r\costh}\]
            よって
            \begin{align*}
                h_r &\coloneq \norm{\partDeriv{\vr}{r}{}} = \norm{\dthv{\sinth\cosph}{\sinth\sinph}{\costh}} = 1\\
                \hth &\coloneq \norm{\partDeriv{\vr}{\theta}{}} = \norm{\dthv{r\costh\cosph}{r\costh\sinph}{-r\sinth}} = r\\
                \hph &\coloneq \norm{\partDeriv{\vr}{\varphi}{}} = \norm{\dthv{-r\sinth\sinph}{r\sinth\cosph}{0}} = r\sinth
            \end{align*}
        \end{proof}
    \section{単位ベクトルに関する3次元Euclid空間との関係}
        \begin{shadebox}
            \begin{align*}
                \ir &= \vix\sinth\cosph + \viy\sinth\sinph + \viz\costh\\
                \ith &= \vix\costh\cosph + \viy\costh\sinph - \viz\sinth\\
                \iph &= -\vix\sinph + \viy\cosph
            \end{align*}
            \begin{align*}
                \vix &= \ir\sinth\cosph + \ith\costh\cosph -\iph\sinph\\
                \viy &= \ir\sinth\sinph + \ith\costh\sinph + \iph\cosph\\
                \viz &= \ir\costh - \ith\sinth
            \end{align*}
        \end{shadebox}
        \noindent(手間が大きい証明)
        \begin{proof}
            \quad\par
            位置ベクトルは
            \[\vr = \dthv{r\sinth\cosph}{r\sinth\sinph}{r\costh}\]
            よって
            \begin{align}
                \ir &\coloneq \frac{1}{h_r}\partDeriv{\vr}{r}{} = \dthv{\sinth\cosph}{\sinth\sinph}{\costh} = \vix\sinth\cosph + \viy\sinth\sinph + \viz\costh \\
                \ith &\coloneq \frac{1}{\hth}\partDeriv{\vr}{\theta}{} = \frac{1}{r}\dthv{r\costh\cosph}{r\costh\sinph}{-r\sinth} = \vix\costh\cosph + \viy\costh\sinph - \viz\sinth \\
                \iph &\coloneq \frac{1}{\hph}\partDeriv{\vr}{\varphi}{} = \frac{1}{r\sinth}\dthv{-r\sinth\sinph}{r\sinth\cosph}{0} = -\vix\sinph + \viy\cosph
            \end{align}
            これを逆に解いてみる。\\
            $(1)\times\costh-(2)\times\sinth$より$\viz = \ir\costh - \ith\sinth$\\
            $(1)\times\sinth+(2)\times\costh$より
            \begin{equation}
                \ir\sinth+\ith\costh = \vix\cosph+\viy\sinph
            \end{equation}
            $(4)\times\cosph-(3)\times\sinph$より$\vix = \ir\sinth\cosph + \ith\costh\cosph -\iph\sinph$\\
            $(3)\times\cosph+(4)\times\sinph$より$\viy = \ir\sinth\sinph + \ith\costh\sinph + \iph\cosph$
        \end{proof}
        \noindent(手間が少ない証明)
        \begin{proof}
            \quad\par
            $\ir,\ith,\iph$の導出までは上の証明と同じ。
            $\ir,\ith,\iph$が直交系を成すことを利用し，$\vix = (\vix\cdot\ir)\ir + (\vix\cdot\ith)\ith + (\vix\cdot\iph)\iph,\quad \viy = (\viy\cdot\ir)\ir + (\viy\cdot\ith)\ith + (\viy\cdot\iph)\iph,\quad \viz = (\viz\cdot\ir)\ir + (\viz\cdot\ith)\ith + (\viz\cdot\iph)\iph$として導出できる。
        \end{proof}
    % \section{Laplacian}
    %     \begin{shadebox}
    %         球座標系に於けるLaplace演算子は次式である。
    %         \[ \Delta = \frac{1}{r^2}\partDerivLong{}{r}{}r^2\partDerivLong{}{r}{} + \frac{1}{r^2}\partDerivLong{}{\theta}{2} + \frac{\cos\theta}{r^2\sin\theta}\partDerivLong{}{\theta}{} + \frac{1}{r^2\sin^2\theta}\partDerivLong{}{\varphi}{2} \]
    %     \end{shadebox}
    %     \begin{proof}
    %         \quad\par
    %         直角座標を引数に取る関数$\phi:\realNumbers^3\to\complexNumbers$は$\mathrm{C}^2$級であるとする。
    %         球座標から直角座標への変換関数を$f:[r,\theta,\varphi]^\top\in\realNumbers^3\mapsto[r\sin\theta\cos\varphi,r\sin\theta\sin\varphi,r\cos\theta]^\top\in\realNumbers^3$とする。
    %         $\hat{\phi}(\bm{u}) := \phi(f(\bm{u}))$とする。
    %         $\vr := [x,y,z]^\top\in\realNumbers^3$とする。
    %         \begin{align*}
    %             \Delta_\vr\phi(\vr) &= \partDeriv{\phi}{x}{2} + \partDeriv{\phi}{y}{2} + \partDeriv{\phi}{z}{2} \\
    %             \begin{split}
    %                 &= \partDerivLong{}{x}{}\parens*{\partDeriv{\hat{\phi}}{r}{}\partDeriv{r}{x}{} + \partDeriv{\hat{\phi}}{\theta}{}\partDeriv{\theta}{x}{} + \partDeriv{\hat{\phi}}{\varphi}{}\partDeriv{\varphi}{x}{}} + \partDerivLong{}{y}{}\parens*{\partDeriv{\hat{\phi}}{r}{}\partDeriv{r}{y}{} + \partDeriv{\hat{\phi}}{\theta}{}\partDeriv{\theta}{y}{} + \partDeriv{\hat{\phi}}{\varphi}{}\partDeriv{\varphi}{y}{}} \\
    %                 &\phantom{=} + \partDerivLong{}{z}{}\parens*{\partDeriv{\hat{\phi}}{r}{}\partDeriv{r}{z}{} + \partDeriv{\hat{\phi}}{\theta}{}\partDeriv{\theta}{z}{} + \partDeriv{\hat{\phi}}{\varphi}{}\partDeriv{\varphi}{z}{}}
    %             \end{split} \tag{1}
    %         \end{align*}
    %         右辺第1項を展開すると次式を得る。
    %         \begin{align*}
    %             &\phantom{=} \partDeriv{\phi}{r}{2}\parens*{\partDeriv{r}{x}{}}^2 + \partDeriv{\hat{\phi}}{r}{}\partDeriv{r}{x}{2} + \partDeriv{\phi}{\theta}{2}\parens*{\partDeriv{\theta}{x}{}}^2 + \partDeriv{\hat{\phi}}{\theta}{}\partDeriv{\theta}{x}{2} + \partDeriv{\phi}{\varphi}{2}\parens*{\partDeriv{\varphi}{x}{}}^2 + \partDeriv{\hat{\phi}}{\varphi}{}\partDeriv{\varphi}{x}{2} \\
    %             &\phantom{=} + 2\partDerivIIHetero{\hat{\phi}}{\theta}{r}\partDeriv{\theta}{x}{}\partDeriv{r}{x}{} + 2\partDerivIIHetero{\hat{\phi}}{\varphi}{\theta}\partDeriv{\varphi}{x}{}\partDeriv{\theta}{x}{} + 2\partDerivIIHetero{\hat{\phi}}{r}{\varphi}\partDeriv{r}{x}{}\partDeriv{\varphi}{x}{}
    %         \end{align*}
    %         式(1)の右辺第2,3項はそれぞれ上式の$x$を$y,z$で置き換えたものになる。
    %         以上より式(1)は次式になる。
    %         \begin{align*}
    %             (1) &= \partDeriv{\hat{\phi}}{r}{2}\bracks*{\parens*{\partDeriv{r}{x}{}}^2 + \parens*{\partDeriv{r}{y}{}}^2 + \parens*{\partDeriv{r}{z}{}}^2} + \partDeriv{\hat{\phi}}{r}{}\parens*{\partDeriv{r}{x}{2} + \partDeriv{r}{y}{2} + \partDeriv{r}{z}{2}} \tag{2} \\
    %             &\phantom{=} + \partDeriv{\hat{\phi}}{\theta}{2}\bracks*{\parens*{\partDeriv{\theta}{x}{}}^2 + \parens*{\partDeriv{\theta}{y}{}}^2 + \parens*{\partDeriv{\theta}{z}{}}^2} + \partDeriv{\hat{\phi}}{\theta}{}\parens*{\partDeriv{\theta}{x}{2} + \partDeriv{\theta}{y}{2} + \partDeriv{\theta}{z}{2}} \tag{3} \\
    %             &\phantom{=} + \partDeriv{\hat{\phi}}{\varphi}{2}\bracks*{\parens*{\partDeriv{\varphi}{x}{}}^2 + \parens*{\partDeriv{\varphi}{y}{}}^2 + \parens*{\partDeriv{\varphi}{z}{}}^2} + \partDeriv{\hat{\phi}}{\varphi}{}\parens*{\partDeriv{\varphi}{x}{2} + \partDeriv{\varphi}{y}{2} + \partDeriv{\varphi}{z}{2}} \tag{4} \\
    %             &\phantom{=} + 2\partDerivIIHetero{\hat{\phi}}{\theta}{r}\parens*{\partDeriv{\theta}{x}{}\partDeriv{r}{x}{} + \partDeriv{\theta}{y}{}\partDeriv{r}{y}{} + \partDeriv{\theta}{z}{}\partDeriv{r}{z}{}} \tag{5} \\
    %             &\phantom{=} + 2\partDerivIIHetero{\hat{\phi}}{\varphi}{\theta}\parens*{\partDeriv{\varphi}{x}{}\partDeriv{\theta}{x}{} + \partDeriv{\varphi}{y}{}\partDeriv{\theta}{y}{} + \partDeriv{\varphi}{z}{}\partDeriv{\theta}{z}{}} \tag{6} \\
    %             &\phantom{=} + 2\partDerivIIHetero{\hat{\phi}}{r}{\varphi}\parens*{\partDeriv{r}{x}{}\partDeriv{\varphi}{x}{} + \partDeriv{r}{y}{}\partDeriv{\varphi}{y}{} + \partDeriv{r}{z}{}\partDeriv{\varphi}{z}{}} \tag{7}
    %         \end{align*}
    %         まず式(5),(6),(7)が0であることを示す。
    %         $J$を$f$のJacobi行列とすると，式(5)の括弧の内側は$J^{-1}$の1行目と2行目の内積である。
    %         球座標系は直交座標系であるため$J$の列は直交系を成す。
    %         これと\ref{行列の列(行)の直交性と逆行列の行(列)の直交性}より$J^{-1}$の行は直交系を成すから，式(5)の括弧の内側は0である。
    %         同様に式(6),(7)も0である。
    %         \par
    %         式(2),(3),(4)を評価する。
    %         $r = \sqrt{x^2+y^2+z^2},\theta = \acos\parens*{z/r} = \acos\parens*{z/\sqrt{x^2+y^2+z^2}},\phi = \atanEx{x}{y}$
    %         であるから次式が成り立つ。
    %         \begin{gather*}
    %             \partDeriv{r}{x}{} = \sin\theta\cos\varphi, \quad \partDeriv{r}{y}{} = \sin\theta\sin\varphi, \quad \partDeriv{r}{z}{} = \cos\theta \\
    %             \partDeriv{\theta}{x}{} = \frac{1}{r}\cos\theta\cos\varphi, \quad \partDeriv{\theta}{y}{} = \frac{1}{r}\cos\theta\sin\varphi, \quad \partDeriv{\theta}{z}{} = -\frac{1}{r}\sin\theta \\
    %             \partDeriv{\varphi}{x}{} = -\frac{\sin\varphi}{r\sin\theta}, \quad \partDeriv{\varphi}{y}{} = \frac{\cos\varphi}{r\sin\theta}, \quad \partDeriv{\varphi}{z}{} = 0 \\
    %             \partDeriv{r}{x}{2} = \frac{1}{r}\parens*{\cos^2\theta\cos^2\varphi + \sin^2\varphi}, \quad \partDeriv{r}{y}{2} = \frac{1}{r}\parens*{\cos^2\theta\sin^2\varphi + \cos^2\varphi}, \quad \partDeriv{r}{z}{2} = \frac{1}{r}\sin^2\theta \\
    %             \partDeriv{\theta}{x}{2} = \frac{1}{r^2\sin\theta}\parens*{-2\sin^2\theta\cos\theta\cos^2\varphi + \cos\theta\sin^2\varphi} \\
    %             \partDeriv{\theta}{y}{2} = \frac{1}{r^2\sin\theta}\parens*{-2\sin^2\theta\cos\theta\sin^2\varphi + \cos\theta\cos^2\varphi} \\
    %             \partDeriv{\theta}{z}{2} = \frac{2}{r^2}\sin\theta\cos\theta \\
    %             \partDeriv{\varphi}{x}{2} = \frac{2\sin\varphi\cos\varphi}{r^2\sin^2\theta}, \quad \partDeriv{\varphi}{y}{2} = -\frac{2\sin\varphi\cos\varphi}{r^2\sin^2\theta}
    %         \end{gather*}
    %         これらを用いて
    %         \begin{align*}
    %             (2) &= \partDeriv{\hat{\phi}}{r}{2} + \frac{2}{r}\partDeriv{\hat{\phi}}{r}{} = \frac{1}{r^2}\partDerivLong{}{r}{}r^2\partDeriv{\hat{\phi}}{r}{} \\
    %             (3) &= \frac{1}{r^2}\partDeriv{\hat{\phi}}{\theta}{2} + \frac{\cos\theta}{r^2\sin\theta}\partDeriv{\hat{\phi}}{\theta}{} \\
    %             (4) &= \frac{1}{r^2\sin^2\theta}\partDeriv{\hat{\phi}}{\varphi}{2}
    %         \end{align*}
    %     \end{proof}